% Round 18. ci_global branch. Independent review required before activation.
\section{Complete height families of non-hypersurface fourfolds}
\label{r18cg:section}

This section constructs a fixed-number-field, Zariski-dense family
of smooth complete intersections of type $(2,5)$ whose full local
strong-G domains occur at every finite place.  A specified prime
has the deep degeneration of CA18, with $\operatorname{rank}N^4=1$.
At the other primes the reductions are good or ordinary quadratic.
The fixed-member Lefschetz argument is proved on the actual quadric
ambient variety; GP17's projective-space statement is not applied
outside its hypotheses.

All local conclusions retain the original valued field, its absolute
value, the precise support of the specified cycle, and every compact
ordinary actual test.  The full-height implication is the frozen
PH16 theorem, with GA17's admissible-lift conclusion on that proved
domain.  No claim about a general variety or about historical
priority for arithmetic pencil constructions is made.

\subsection{A smooth integral quadric and its section space}

Fix any preassigned number field $F_{\mathrm{base}}$.  Let
\[
 \mathcal Y=V(q_2)\subset\mathbf P^6_{\mathbf Z},\qquad
 q_2=X_0X_1+X_2X_3+X_4X_5+X_6^2.
 \tag{CG18.1}\label{r18cg:quadric}
\]
This is a smooth projective relative quadric of dimension five
over $\mathbf Z$.  Indeed, in every characteristic a common zero
of its six partial derivatives with respect to $X_0,\ldots,X_5$
has these six coordinates zero.  The equation then forces
$X_6=0$, which is not a projective point.  This also proves
smoothness in characteristic two.  Every geometric fiber is
connected, hence integral: for example the standard projective
hypersurface sequence gives $H^0(\mathcal Y_{\bar k},\mathcal O)=\bar k$,
and smooth components of a connected scheme cannot meet or be
disjoint components of it.

Let $V_{\mathbf Z}$ be the free module of degree-five homogeneous
forms in the seven displayed variables.  We work with actual forms
in this space and their restrictions to $\mathcal Y$.
They supply the complete section space of $\mathcal O_{\mathcal Y}(5)$:
the sequence
\[
 0\longrightarrow\mathcal O_{\mathbf P^6}(3)
    \xrightarrow{q_2}\mathcal O_{\mathbf P^6}(5)
    \longrightarrow\mathcal O_{\mathcal Y}(5)\longrightarrow0
 \tag{CG18.2}\label{r18cg:section-sequence}
\]
and the vanishing of the intermediate cohomology of projective
space give this surjectivity over any field.  One can alternatively
keep the full ambient coefficient space throughout; its linear
kernel consists of the multiples $q_2V_{\mathbf Z,3}$ and does not
affect any of the arguments below.

\begin{lemma}[Smooth sections and isolated-node integrality]
\label{r18cg:smooth-integral}
Over every field $k$, a nonempty geometrically dense open set in
$\mathbf P(V_k)$ consists of sections cutting a smooth fourfold
on $Y_k$.  A section whose divisor on $Y_{\bar k}$ has at most
isolated singular points defines a geometrically integral fourfold.
\end{lemma}

\begin{proof}
Degree-five forms give every first jet at each point of
$\mathbf P^6_{\bar k}$: at $[1:0:\cdots:0]$ use $X_0^5$ and
$X_0^4X_i$ for $1\le i\le6$.  Taking the quotient to the
first jets on the smooth fivefold $Y$ preserves surjectivity.
The incidence of a section with a marked singular point is
therefore a projective bundle over $Y$ of dimension one less
than $\mathbf P(V_k)$.  Its image is closed, since $Y$ is
proper, and is a proper subset by that dimension bound.
Its complement is the asserted geometrically nonempty open.
The argument applies also in small positive characteristic.

For integrality, a nonzero section defines a Cartier divisor
on the smooth integral $Y_{\bar k}$.  Its local rings are
Cohen--Macaulay, hence satisfy Serre's condition $S_2$.
An isolated singular locus in dimension four gives $R_1$.
Thus the divisor is normal.  It is connected: the exact sequence
$0\to\mathcal O_Y(-5)\to\mathcal O_Y\to\mathcal O_X\to0$
gives $H^0(X,\mathcal O_X)=\bar k$, because
\[
                   H^1(Y,\mathcal O_Y(-5))=0.
\]
For a direct check of this vanishing, use
$0\to\mathcal O_{\mathbf P^6}(-7)\to
\mathcal O_{\mathbf P^6}(-5)\to\mathcal O_Y(-5)\to0$
and the intermediate projective-space cohomology vanishings.
A connected normal scheme has just one irreducible component,
since normal irreducible components are open and closed.
It is therefore integral.  This proves the geometric assertion;
in particular it applies to the smooth members and to the
ordinary-node members used below.
\end{proof}

\subsection{Lefschetz pencils through a fixed smooth section}

\begin{lemma}[The fixed-member open on the quadric]
\label{r18cg:fixed-member}
Let $k$ have characteristic zero, and let $F_0\in V_k$ cut a
smooth fourfold on $Y_k$.  There is a nonempty Zariski open
subset of $V_k$ of forms $F_1$ for which:
\begin{enumerate}
\item the section $F_1=0$ on $Y_k$ is smooth;
\item the base $F_0=F_1=0$ on $Y_k$ is smooth of dimension three;
\item every singular geometric member of the pencil
      $uF_0+vF_1=0$ on $Y_k$ has exactly one ordinary quadratic
      singularity, and its parameter line meets the section
      discriminant transversely there.
\end{enumerate}
\end{lemma}

\begin{proof}
We prove the assertion geometrically.  Write
$\mathbf P(V_k)=\mathbf P^M$ and let
\[
 I=\{([f],x):j_x^1(f|_Y)=0\}\subset\mathbf P^M\times Y.
 \tag{CG18.3}\label{r18cg:incidence}
\]
First-jet surjectivity makes this a smooth irreducible projective
bundle of dimension $M-1$.  The quadratic term at a marked point
can be prescribed arbitrarily: degree-five forms give every
second jet on projective space, using the monomials
$X_0^5$, $X_0^4X_i$ and $X_0^3X_iX_j$, and their images give
every second jet on $Y$.  Hence degeneracy of that quadratic
term is a proper closed subset of $I$, of dimension at most
$M-2$.

First-jet evaluations at two distinct points are jointly
surjective.  Move the two points in projective space to its
first two coordinate points.  The monomials prescribing the
two first jets are $X_0^5,X_0^4X_i$ and
$X_1^5,X_1^4X_j$ with the respective remaining indices.
They are distinct and have zero first jet at the opposite
point.  Restricting to the two smooth tangent spaces of $Y$
preserves this surjectivity.  Thus the incidence with two
distinct marked singular points has dimension
\[
          2\dim Y+M-2(\dim Y+1)=M-2.
\]
Its image closure in the coefficient space has dimension at
most $M-2$.

First forget only the second marked point in the two-point
incidence.  Its image closure in $I$ has dimension at most
$M-2$.  Removing that subset and the degenerate-quadratic
subset leaves a nonempty open of $I$ parametrizing sections
with one singular point and invertible Hessian.  The projection
to coefficients has singleton fibers on this open, so its image
has dimension $M-1$.  The full image of $I$ is consequently an
irreducible discriminant hypersurface $\Delta$.
The images in coefficient space of the two bad incidences
have closures of dimension at most $M-2$.  Their removal
gives a dense open of $\Delta$ consisting of one-node sections.
At such a point the critical-point equations can be solved
etale locally as coefficients vary, because their Hessian is
invertible.  The remaining critical-value equation has a
nonzero differential, namely evaluation of the coefficient
variation at the marked point.  The discriminant is smooth
there.  Denote this one-node smooth open by $\Delta^\circ$;
its complement in $\Delta$ has dimension at most $M-2$.

Projection from $[F_0]\notin\Delta$ gives a finite morphism
\[
                         \Delta\longrightarrow\mathbf P^{M-1}
\]
to the space of lines through $[F_0]$.  Its fibers are finite
line intersections, and properness turns quasi-finiteness into
finiteness.  The map is generically etale in characteristic
zero.  Exclude its branch locus and the image of
$\Delta\setminus\Delta^\circ$.  What remains is a nonempty
open set of directions whose lines have precisely the required
transverse one-node intersections.

There is another nonempty open set of directions with smooth
base on the fixed $X_0=Y\cap V(F_0)$.  The restrictions of
degree-five forms give a very ample, base-point-free system
on this smooth fourfold, so the usual characteristic-zero
Bertini theorem applies to it.  Intersect the two open sets
of directions.  Lift to choices of $F_1$ and exclude the
discriminant itself.  This gives the asserted nonempty open.
At each base point the differentials of $F_0,F_1$ on $Y$
are independent, so every pencil member is smooth there.
All constructions descend to $k$.
\end{proof}

\paragraph{Classical inputs.}
The Bertini and Lefschetz-pencil setting is classical; see
Jannsen--Saito, \emph{Bertini theorems and Lefschetz pencils over
discrete valuation rings}, \url{https://arxiv.org/pdf/0911.1470},
Sections 0--2.  The proof above records the fixed-member and
degree-five jet assertions needed here, instead of assuming
that an unspecified generic pencil can be arranged through
the already chosen section.

\subsection{One fixed number field handles the small residue fields}

Choose an integral $F_0\in V_{\mathbf Z}$ cutting a smooth
fourfold over $\mathbf Q$; the preceding nonempty open and
density of rational points ensure that such a choice exists.
Let $S_0$ be a finite set of rational primes containing $2$
and every prime where $F_0=0$ on $\mathcal Y$ does not have
smooth four-dimensional reduction.  Its finiteness follows
by spreading this fixed proper smooth generic member.

For each $p\in S_0$, Lemma~\ref{r18cg:smooth-integral} gives
a nonempty geometric open of smooth quintic sections on
$Y_{\mathbf F_p}$.  Choose one point of this open over a
finite field $\mathbf F_{p^{f_p}}$.  Such a point exists
because each of its finitely many algebraic coordinates is
defined over a finite extension of $\mathbf F_p$.
Choose an integer $N\ge2$ divisible by every $f_p$ and view
all these chosen sections over $\mathbf F_{p^N}$.
There is no assertion that a fixed-degree smooth section
exists already over each original small residue field.

For each $p\in S_0$ choose a monic irreducible polynomial of
degree $N$ over $\mathbf F_p$.  The coefficient Chinese
remainder theorem lifts all these polynomials to one monic
$g\in\mathbf Z[T]$ of degree $N$.
It is irreducible over $\mathbf Q$ by reduction at any of
these primes.  Put $L=\mathbf Q[T]/(g)$ and
\[
                              F=F_{\mathrm{base}}L.
 \tag{CG18.4}\label{r18cg:field}
\]
At each $p\in S_0$, the ring $\mathbf Z_p[T]/(g)$ is the
unramified degree-$N$ discrete valuation ring, since its
reduction is irreducible and separable.  Thus $p$ is inert
and unramified in $L$, and its residue field is
$\mathbf F_{p^N}$.  This also follows from the fact that
the discriminant of $g$ is a unit at $p$.
The compositum $F$ need not be unramified at these primes;
its every residue field above $p$ contains the indicated
residue field of $L$, which is all that is needed.

Lift the chosen smooth section over $\mathbf F_{p^N}$ to
an actual degree-five form $G_p\in\mathcal O_L[X_0,\ldots,X_6]$,
and regard it also over $\mathcal O_F$.
At every prime of $F$ above $p$, its reduction remains a
smooth section by scalar extension.
The field $F$ and these finitely many forms are now fixed,
before selecting any pencil or its parameters.

\subsection{A fixed deep local condition and the coefficient lattice}

Over $\mathbf Q$ choose five linear forms
$\ell_1,\ldots,\ell_5$ whose hyperplane sections on $Y$
have all the geometric SNC and complete-intersection
properties in CA18.  Their fivefold intersection consists
of two geometric points.  Choose a quintic $Q$ cutting
every positive-dimensional intersection transversely and
avoiding those two points.  These are finitely many
nonempty characteristic-zero Bertini open conditions.
Clear denominators and put $P=\prod_{i=1}^5\ell_i$.

Choose a rational prime $q$ outside $S_0$, outside the
discriminant of the already fixed field $F$, and outside
the finite set where these transverse intersection data
fail to spread.  The prime can avoid any additional
prescribed finite set.  In particular $q$ is odd and is
a uniformizer at every completion $F_w$ with $w\mid q$.

Set
\[
                 B=\prod_{p\in S_0}p,\qquad M=q^2B.
 \tag{CG18.5}\label{r18cg:MB}
\]
Choose a degree-five form $F_1$ with coefficients in
$\mathcal O_F$ satisfying
\[
 \begin{aligned}
 F_1&\equiv P+qQ&&\pmod{q^2\mathcal O_F},\\
 F_1&\equiv G_p&&\pmod{p\mathcal O_F}\quad(p\in S_0).
 \end{aligned}
 \tag{CG18.6}\label{r18cg:coefficient-CRT}
\]
These ideals are pairwise coprime, so the solution set is
one nonempty full affine coefficient lattice, a translate
of $M\mathcal O_F^{\dim V}$.
It is Zariski dense over $F$: it contains a product of
infinite one-variable arithmetic progressions, and the
one-variable polynomial root bound, applied successively,
proves that no nonzero polynomial vanishes on it.
Its intersection with the fixed-$F_0$ open of
Lemma~\ref{r18cg:fixed-member} is therefore nonempty and
Zariski dense.  Restrict $F_1$ to that intersection.

After fixing this $F_1$, spread its pencil over
$\mathcal O_F[1/S]$ for a finite set of prime ideals $S$.
We choose $S$ to contain all primes over $q$ and $S_0$.
Outside $S$, the base and total pencil are smooth, the
critical scheme is finite etale and maps as a closed
immersion into the parameter line, and its Hessians are
invertible.  We also require that the two sections are
linearly independent in each residue field.

Here is why that spread is available.  Over $F$ the
incidence divisor $uF_0+vF_1=0$ in $Y\times\mathbf P^1$
is smooth: away from the base a parameter derivative is
nonzero, and at the base the two section differentials
are independent.  Every fiber is a nonzero Cartier
divisor, so the pencil is flat.  The lemma gives finitely
many nondegenerate critical points with distinct geometric
values; their scheme is finite etale in characteristic
zero and has a closed-immersion map to $\mathbf P^1$.
All data are of finite presentation.  After removing
finitely many prime ideals, smoothness, finite etaleness,
closed immersion and the invertible Hessians spread.
One first removes any vertical components of the critical
locus; no extra critical points then remain on the spread.
The independence condition excludes zero sections as fibers.
Every geometric member outside $S$ is consequently smooth
or has exactly one ordinary quadratic singularity.
No uniform set $S$ is asserted as $F_1$ varies.

\subsection{The arithmetic members and all finite places}

Let $T$ be the finite set of rational primes below the
prime ideals of $S$.  Choose integers $a$ in the infinite
arithmetic progression
\[
 \begin{aligned}
 a&\equiv1\pmod{q^2},&
 a&\equiv1\pmod B,\\
 a&\equiv0\pmod p&&
       (p\in T\setminus(S_0\cup\{q\})).
 \end{aligned}
 \tag{CG18.7}\label{r18cg:parameter-CRT}
\]
All distinct moduli are coprime.  Discard the finitely
many parameters whose characteristic-zero pencil member
is singular, and put
\[
    X_a=Y_F\cap V(MF_0+aF_1)\subset\mathbf P^6_F.
 \tag{CG18.8}\label{r18cg:family}
\]
Every remaining member is a smooth geometrically integral
complete-intersection fourfold of type $(2,5)$.

\begin{theorem}[Full local G and complete heights for a non-hypersurface family]
\label{r18cg:height-family}
For the fixed number field $F$ of \eqref{r18cg:field},
every member \eqref{r18cg:family} has the following properties.
\begin{enumerate}
\item At every finite place $w$ of $F$, every specified
      geometrically homologically trivial rational cycle on
      $(X_a)_{F_w}$ satisfies strong G in every codimension,
      with its exact original support and every original
      compact ordinary actual test.
\item At every $w\mid q$ it has the actual reduced strictly
      semistable model of CA18 and
      $\operatorname{rank}N^4=1$.  It has no potential good
      reduction there and no strict semistable dominant
      witness of crossing depth less than four.
\item The preceding full local assertion holds for all new
      cycles after every finite extension $F'/F$ unramified
      at all places above $q$.  At the other places no
      ramification restriction is imposed.
\item For every such $F'$ (in particular $F'=F$) and every
      $1\le r\le4$, the frozen PH16 theorem gives the
      canonical real-valued Beilinson--Bloch pairing on
      the complete displayed homological Chow groups:
      \[
      \operatorname{CH}^r((X_a)_{F'})_{\mathbf Q,\mathrm{geom\,hom}}
       \ \times\
      \operatorname{CH}^{5-r}((X_a)_{F'})_{\mathbf Q,\mathrm{geom\,hom}}
                        \longrightarrow\mathbf R.
      \tag{CG18.9}\label{r18cg:height}
      \]
\end{enumerate}
Every geometric $X_a$ is not isomorphic to any smooth
projective hypersurface fourfold, even after a field
extension.  As $F_1$ varies through the fixed lattice
\eqref{r18cg:coefficient-CRT} in the fixed-member open
and $a$ varies through its permitted progression, the
resulting members are Zariski dense in the space of
smooth sections of $\mathcal O_{Y_F}(5)$.
The field $F$, the quadric, and the distinguished prime
are fixed throughout this density statement.
\end{theorem}

\begin{proof}
At a place $w$ above $q$, the exact integral equation is
\[
 MF_0+aF_1=P+qQ_a,\qquad
 Q_a=\frac{MF_0+aF_1-P}{q},\qquad \bar Q_a=\bar Q.
 \tag{CG18.10}\label{r18cg:local-product}
\]
Coefficientwise integrality and the last equality follow
from $q^2\mid M$, \eqref{r18cg:coefficient-CRT}, and
$a\equiv1\pmod{q^2}$.  Because $q$ is unramified in $F$,
it has uniformizer exponent one in the native ring
$\mathcal O_{F_w}$.  All special-fiber transverse factor
data hold after the residue extension.  CA18 therefore
applies to the entire native cycle group over $F_w$,
giving full strong G and the rank-one fourth monodromy.

If $w$ lies above $p\in S_0$, the coefficient $M$ reduces
to zero and $a$ to one.  The equation reduces to the
smooth section $G_p$ on the special quadric.  This is
good reduction, including at every place of characteristic
two.  If $w$ lies above a prime
$p\in T\setminus(S_0\cup\{q\})$, then $M$ is a unit
and $a$ reduces to zero.  The reduction is the smooth
section $F_0$.  This also gives good reduction.  Although
$S$ was a set of prime ideals, imposing $a\equiv0\pmod p$
does no harm at the other places above the same rational
prime: $F_0$ is smooth at all of them by the definition
of $S_0$.

At every remaining place $w$, the pencil is in its
spread good set and $M$ is a unit.  The parameter
$[M:a]$ has a well-defined integral specialization.
Its special member is smooth or has one ordinary
quadratic point.  It is geometrically integral by
Lemma~\ref{r18cg:smooth-integral}.  In all cases the
equation gives a flat projective model: it is a
nonzero Cartier equation also after reduction, so
the divisor sequence has no uniformizer torsion.
Its integral smooth generic fiber then makes the
flat model integral as well.

The remaining residue characteristics are odd.
The frozen ON17 theorem applies to this flat integral
model of ordinary quadratic reduction.  Its hypotheses
do not require a projective-space hypersurface: ON17
constructs the native smooth local ambient from the
flat ordinary quadratic local ring.  Here the smooth
quadric already supplies such an ambient.
The possible finite node thickness is arbitrary and
is retained by ON17 and QF17.  Those theorems therefore
give full strong G in every degree at these places.
At the good places use the frozen good-reduction theorem.
This proves the full placewise statement.

For a finite $F'/F$ unramified above $q$, the exact
product model \eqref{r18cg:local-product} still has
uniformizer exponent one at each place above $q$.
Apply CA18 afresh to every specified new cycle over
each new completion; none is presumed to descend as
a pullback.  At good places, base change preserves
good reduction.  At ordinary quadratic places,
ON17's finite-extension corollary applies anew;
the thickness is multiplied by the ramification
index, and all its native auxiliary degrees and
absolute values remain those of that theorem.
This proves the asserted extension domain.

Every finite place and both complementary cycle
codimensions now satisfy precisely the hypotheses
of PH16.  Its full-domain height implication gives
\eqref{r18cg:height}; GA17 supplies the admissible
regular-model lifts on this established domain.
No unproved general-height premise is used.

The rank-one monodromy and crossing-depth assertions
are CA18 with ambient degree two, five linear factors
and the frozen NM17 obstruction.  The non-hypersurface
assertion is intrinsic: adjunction gives $K_{X_a}=0$;
a hypothetical hypersurface fourfold isomorphism
would force degree six, whereas its top Chern degree
would be $2610$ instead of the actual $2190$, as
computed directly in CA18.  The argument does not
presume any polarization compatibility.

For density, let a polynomial $b$ in the affine form
coefficients over $F$ vanish on every constructed
$MF_0+aF_1$.  For each allowed $F_1$, it vanishes
at infinitely many integers $a$ on that affine line.
The polynomial $b(MF_0+tF_1)$ is therefore identically
zero.  At $t=1$ this gives $b(MF_0+F_1)=0$ for every
allowed $F_1$.  Those forms are Zariski dense in the
full affine coefficient space, so translation gives
$b=0$.  The linear surjection to
$H^0(Y_F,\mathcal O(5))$, followed by projectivization,
proves density in the section parameter space and
then in its smooth open.  The set $S$ and the
progression can depend on $F_1$; infinitely many
parameters for each $F_1$ suffice for the argument.
\end{proof}

\paragraph{Additional finite-extension conclusion, conditional only on RI18.}
If RI18's inertia-and-norm theorem is independently activated,
then the unramified restriction at $q$ in
Theorem~\ref{r18cg:height-family} can be removed.
Indeed CA18 gives full G anew after every unramified
extension of each native $F_w$, exactly the hypothesis
needed by RI18.  RI18 then covers all cycles over every
finite extension of $F_w$, using its full-degree actual
trace descent.  All other places already have that
extension stability by good reduction and ON17.
PH16 therefore gives \eqref{r18cg:height} over every
finite $F'/F$.  This paragraph does not count RI18 as
proved before its separate independent review is complete.

\paragraph{Precise scope and comparison with earlier domains.}
For every preassigned $F_{\mathrm{base}}$ the construction
uses a finite extension $F/F_{\mathrm{base}}$ and then
gives density over that one fixed $F$.  It does not
claim that the displayed family descends to
$F_{\mathrm{base}}$, nor that every smooth type-$(2,5)$
fourfold belongs to it.  The geometric fibers are
not hypersurfaces; their deep degeneration is not
a new model of a potentially good generic fiber.
Thus this is an explicit complete-height class
beyond the hypersurface families GP17 and GF18.
No classification of overlap with every earlier
motivic sector is asserted.  Zariski density supplies
no deformation or specialization principle proving
G for all complete intersections.  General strong G
and arbitrary-variety full unconditional heights
remain open.
